\documentclass[a4paper,12pt]{article}
\usepackage[utf8]{inputenc}
\usepackage[T1]{fontenc}
\usepackage[french]{babel}
\usepackage{geometry}
\geometry{margin=2.5cm}

\title{Compte rendu texte 1}
\author{Mignard Mael}
\date{\today}

\begin{document}

\maketitle

Ce document, publié dans le Guardian en septembre 2021 , explique que les entreprises américaines refusent de voir à quel point leurs employés ont changé après la pandémie.\\ \\
Le texte rappelle que le Covid-19 a détruit le "rêve américain". 
En très peu de temps, 22 millions de personnes ont perdu leur travail et beaucoup sont tombées dans la pauvreté. 
Cette crise a montré que la plupart des gens n'avaient pas assez d'argent de côté pour survivre plus de trois mois sans salaire. 
En plus de ces difficultés financières, la santé mentale des travailleurs s'est fortement dégradée, avec beaucoup plus de cas de dépression.\\ \\
Aujourd'hui, les patrons veulent reprendre le travail comme avant, en se concentrant uniquement sur la productivité.
Mais les salariés ne sont plus d'accord. 
Beaucoup ont compris que le travail ne doit pas être toute leur vie. 
Ils veulent maintenant plus de temps libre ou un métier qui a vraiment du sens. \\ \\
En conséquence, on voit un nombre record de démissions et de grèves. L'auteur conclut que le système actuel est cassé et qu'il n'est pas prêt pour affronter les prochaines crises mondiales.


\end{document}